\section{Polynomial features}
\label{sec:poly}

\(\mathcal{H}_l\) cannot approximate a general continuous target.
On any closed bounded \(B\subset\R^d\), polynomials can: for every continuous \(g\colon B\to\R\) and every \(\delta>0\) there is a polynomial \(f\) with \(\lvert f(x)-g(x)\rvert<\delta\) for all \(x\in B\) \cite{weierstrass1885,stone1948}.
Write \(\mathcal{H}_p\) for that class.

Every \(f\in\mathcal{H}_p\) is linear in a monomial coordinate.
One variable, every degree:
\begin{equation}
f(x)=\sum_{k\in\N} w_k x^k = \langle w,\tilde\varphi(x)\rangle,
\qquad
\tilde\varphi(x)=(1,x,x^2,x^3,\ldots).
\label{eq:phi-1d}
\end{equation}
In \(d\) variables the same statement uses one index per coordinate,
\begin{equation}
f(x)=\sum_{\alpha\in\N^d} w_\alpha x^\alpha = \langle w,\tilde\varphi(x)\rangle.
\label{eq:phi-d}
\end{equation}
Raw monomials are unusable as features: at \(x=1\) and \(w=(1,1,\ldots)\) the inner product diverges, and for \(d=100\) the monomials of degree \(5\) already number almost \(10^8\).

Scale each monomial by a Gaussian bump and a factorial,
\begin{equation}
\varphi(x)
= e^{-\lVert x\rVert^2/2}
\Biggl[\frac{x^\alpha}{\sqrt{\alpha!}}\Biggr]_{\alpha\in\N^d}.
\label{eq:phi}
\end{equation}
Then \(\lVert\varphi(x)\rVert=1\) for every \(x\), there is no degree cutoff, and the model remains a universal approximator.
The fit is still linear regression on features,
\begin{equation}
f(x)=\langle w,\varphi(x)\rangle.
\label{eq:poly-fit}
\end{equation}
The one-variable map is the same formula with \(\lVert x\rVert^2=x^2\) and \(\alpha=k\).
